\section{Introduction}
\label{sec:intro}

Reasoning-tuned language models answer hard questions by spending more
computation at inference time: they sample long chains of thought, draw several
candidates, and verify or vote among them
\cite{wei2022cot,wang2023selfconsistency,snell2024scaling,deepseek2025r1}. The
gain is real but expensive, and it is wasted on the many requests that a direct
answer would have handled \cite{chen2024overthinking,zhang2025adaptthink}. This
is the engineering form of a very old observation: cognition has a fast,
automatic mode and a slow, effortful one, and competence lies in knowing when to
switch \cite{kahneman2011thinking,evans2013dual,sloman1996two}.

The Islamic intellectual tradition developed a rich vocabulary for the same
observation long before dual-process psychology. The Qur'an speaks of the
\emph{nafs} that ``commands to evil'' (Q~12:53), the \emph{nafs} that reproaches
itself (Q~75:2) and the \emph{nafs} ``at rest'' (Q~89:27)
\cite{AbdelHaleem2008Quran}; al-Ghazali describes the heart (\emph{qalb}) as a
king whose vizier is the intellect (\emph{\textquoteleft aql}) and whose unruly
soldiers are appetite and anger \cite{Ghazali2010Marvels}; and the Sufi
disciplines of \emph{muraqaba} (vigilance) and \emph{muhasaba} (self-accounting)
institutionalise the daily audit of one's own conduct
\cite{Picken2011,Qushayri2007Epistle}. We use this vocabulary as a design
metaphor. We make no claim about the metaphysics of the soul, and we say so
once here and again in Section~\ref{sec:islamic}: every mapping in this paper is
an analogy between a classical \emph{description of a faculty} and an
\emph{engineering component}.

\subsection{Problem}
The project that this report serves has a concrete constraint set: build a
dual-process assistant from pretrained Qwen models, serve it on a CPU-only
machine, fine-tune it with Unsloth, and expose it through a Rust web
application with a chat room, a performance-test console and an evaluation
dashboard. An existing draft design \cite{draft2026} targeted 27B--35B models
on 80\,GB GPUs and reviewed twelve formal results; it did not address CPU
feasibility, the model roster, the serving stack, or a buildable
specification. This report assesses that draft, keeps what is sound, corrects
what is not, and extends it into a mathematically principled and feasible
design.

\subsection{Contributions}
\begin{enumerate}
\item A faithful, sourced account of the classical concepts we borrow, with
the disagreements among scholars noted and the analogy kept separate from
theology (Section~\ref{sec:islamic}).
\item A formal model of a verifier-gated dual-process system with
seventeen stated results, each with explicit assumptions and an engineering
consequence (Section~\ref{sec:math}). The results correct two defects of the
draft: its router rule was hand-tuned rather than derived, and its
self-refinement loop was not shown to be stable.
\item A verified feasibility study of the stack as of October 2026: which Qwen
variants exist below 9B parameters, which run in vLLM-CPU and llama.cpp, what
Unsloth can and cannot do without a GPU, which Rust crates to use, and what a
four-core AVX2 machine actually delivers (Sections~\ref{sec:stack} and
\ref{sec:feasibility}).
\item A component architecture with a termination guarantee, an assurance-level
model for answers, and a self-model that is unbiased under routing
(Section~\ref{sec:arch}).
\item An evaluation plan whose gates are statistically sized, and a summary of
the companion specification and roadmap (Sections~\ref{sec:spec} and
\ref{sec:eval}).
\end{enumerate}

\subsection{What this report does not claim}
No accuracy, calibration or cost result for \Sys\ has been measured. Where a
number appears it is (a) a measurement of an existing model on the reference
machine, clearly marked, (b) a value quoted from a cited paper, or (c) a target
or estimate. The roadmap (Section~\ref{sec:spec}) schedules every hypothesis
for validation.
